\documentclass{article}

\usepackage{amsmath,amssymb}
\usepackage[margin=1in]{geometry}
\usepackage{xurl}
\usepackage[hidelinks]{hyperref}
\setlength{\emergencystretch}{2em}

% Shared commands for notes in docs/paper-gaps/.

\DeclareUnicodeCharacter{2080}{\ifmmode{}_{0}\else\textsubscript{0}\fi}
\DeclareUnicodeCharacter{2081}{\ifmmode{}_{1}\else\textsubscript{1}\fi}
\DeclareUnicodeCharacter{2082}{\ifmmode{}_{2}\else\textsubscript{2}\fi}
\DeclareUnicodeCharacter{2099}{\ifmmode{}_{n}\else\textsubscript{n}\fi}

\newcommand{\Lean}{\textsc{Lean}}
\ExplSyntaxOn
\NewDocumentCommand{\leanid}{m}{
  \ifmmode
    \textsf{\detokenize{#1}}
  \else
    \group_begin:
    \urlstyle{sf}
    \tl_set:Nn \l_tmpa_tl {#1}
    \tl_replace_all:Nnn \l_tmpa_tl {\_} {_}
    \exp_args:NV \path \l_tmpa_tl
    \group_end:
  \fi
}
\ExplSyntaxOff
\providecommand{\tr}{\operatorname{tr}}
\newcommand{\tnleanIssueBase}{https://github.com/LionSR/TNLean/issues/}
\newcommand{\tnleanPrBase}{https://github.com/LionSR/TNLean/pull/}
\newcommand{\tnleanDocBase}{https://sirui-lu.com/TNLean/docs/}
\newcommand{\ghissue}[1]{\href{\tnleanIssueBase#1}{GitHub issue \##1}}
\newcommand{\ghpr}[1]{\href{\tnleanPrBase#1}{PR \##1}}
\newcommand{\doclink}[1]{\href{\tnleanDocBase#1.html}{\path{#1}}}

% Dirac notation and the identity operator, as in blueprint/src/macros/common.tex.
% \providecommand keeps note-local definitions authoritative.
\usepackage{dsfont}
\providecommand{\ket}[1]{|#1\rangle}
\providecommand{\bra}[1]{\langle #1|}
\providecommand{\braket}[2]{\langle #1 | #2 \rangle}
\providecommand{\ketbra}[2]{|#1\rangle\!\langle #2|}
\providecommand{\Id}{\mathds{1}}
\providecommand{\one}{\mathds{1}}
\providecommand{\C}{\mathbb{C}}
\providecommand{\MN}[1]{M_{#1}(\C)}
\providecommand{\MD}{M_D(\C)}

% External referencing for paper sources.  The build helper writes sanitized
% auxiliary files under build/paper-aux/ and excludes duplicated source labels.
\usepackage{xr}

% Import a generated paper auxiliary file with a caller-supplied label prefix.
% The candidate paths support builds from docs/paper-gaps/, the repository root,
% and build/paper-aux/ itself.
\newcommand{\importpaperaux}[2]{%
  \IfFileExists{../../build/paper-aux/#2.aux}{%
    \externaldocument[#1]{../../build/paper-aux/#2}%
  }{%
    \IfFileExists{build/paper-aux/#2.aux}{%
      \externaldocument[#1]{build/paper-aux/#2}%
    }{%
      \IfFileExists{#2.aux}{%
        \externaldocument[#1]{#2}%
      }{}%
    }%
  }%
}

% General external reference macros
\newcommand{\paperref}[2]{\ref{#1-#2}}
\newcommand{\papereqref}[2]{(\ref{#1-#2})}

% CPSV16 source (1606.00608 / MPDO-22-12-17-2.tex) external document and shortcuts
\importpaperaux{cpsv16-}{1606.00608/MPDO-22-12-17-2-tnlean}
\newcommand{\cpsvref}[1]{\paperref{cpsv16}{#1}}
\newcommand{\cpsveqref}[1]{\papereqref{cpsv16}{#1}}

% Verdict marker: \gapnote{<kind>}{<status>}. Kind is the policy
% classification (clarification, local-correction, scope-restriction,
% unfaithful, false-source, open-gap); status is open, wip (elimination
% underway), resolved, or historical. Machine-read by the site index;
% prints a verdict line.
\newcommand{\gapnote}[2]{\par\noindent\textbf{Verdict:} #1 (#2).\par}


\title{The Missing Algebraic Specification of the PEPS Boundary-RFP Diagrams}
\author{}
\date{2026-08-30}

\begin{document}
\maketitle
\gapnote{open-gap}{open}

\section{Source assertion}

Section~4.2 of Cirac--P\'erez-Garc\'ia--Schuch--Verstraete considers a
homogeneous square-lattice PEPS tensor $A$. It assumes the graphical
renormalization relation in Figure~\cpsvref{PEPS-RFP}, where the equality is
understood up to an isometry on the physical indices. The paper then draws a
local tensor for the one-dimensional boundary state and asserts that the PEPS
relation produces the two maps in the following diagram, ``which gives the
desired $T$ and $S$'' of Definition~\cpsvref{RFPMixedTS}
\cite[lines~694--729]{Cirac2016MPDO_arXiv}.
\footnote{Source file:
    \path{Papers/1606.00608/MPDO-22-12-17-2.tex}:694--729.}

Thus the paper-facing conclusion is that the boundary MPO is a
renormalization fixed point: there should be trace-preserving completely
positive maps $\mathcal S$ and $\mathcal T$ such that
\begin{align}
    \mathcal S[M_2(X)] &= M_1(X),
    \label{eq:cpsv16_peps_boundary_rfp_specification_coarse_graining} \\
    \mathcal T[M_1(X)] &= M_2(X)
    \label{eq:cpsv16_peps_boundary_rfp_specification_refinement}
\end{align}
for every virtual operator $X$. The existence and channel properties of these
maps are conclusions of the PEPS renormalization relation, not additional
hypotheses.

\section{The local boundary tensor}

The local boundary tensor in the diagram at source lines~721--724 has two
ket--bra contractions. To make this explicit, write
$A^s_{l,r,u,b}$ for a PEPS tensor whose virtual legs are ordered left, right,
up, and down. Choosing the down leg as the inward leg, the physical leg $s$
and inward virtual leg $b$ are contracted between the ket and bra copies:
\begin{align}
    (E_A)^{u,u'}_{(l,l'),(r,r')}
        &=
    \sum_{s=0}^{d-1}\sum_{b=0}^{D_{\mathrm v}-1}
    A^s_{l,r,u,b}\,
    \overline{A^s_{l',r',u',b}}.
    \label{eq:cpsv16_peps_boundary_rfp_specification_local_tensor}
\end{align}
Consequently, $E_A$ has boundary physical dimension $D_{\mathrm v}$ and MPO
bond dimension $D_{\mathrm h}^2$. Here the bra legs are labelled after
reflecting the bra copy into the ket coordinate frame. This is why the
contracted leg is the down coordinate $b$ in both component functions; in the
unreflected stacked picture it is the ket-down leg joined to the bra-up leg.

This local formula is formalized without a fixed-point or channel
hypothesis.
\footnote{Formal declaration:
    \leanid{TNLean.PEPS.localBoundaryMPOTensor}.}

\section{Data absent from the diagrams}

The source passage does not give an algebraic form of the PEPS blocking
relation. In particular, it does not specify the domain and codomain of the
physical isometry or the identifications between the external virtual spaces
before and after blocking. It also does not specify the spaces, orientation,
choice, or normalization of the left and right transfer fixed points used to
form the boundary state.

The refinement part of the final diagram inserts an unlabelled local factor.
Its space and normalization are not stated. The coarse-graining and refinement
pictures therefore do not by themselves define linear maps on the full
one-site and two-site boundary matrix algebras. An algebraic proof of
Eqs.~\eqref{eq:cpsv16_peps_boundary_rfp_specification_coarse_graining} and
\eqref{eq:cpsv16_peps_boundary_rfp_specification_refinement} must additionally
show that the resulting maps are completely positive and trace-preserving on
those full algebras.

These are not merely missing proof steps for a stated algebraic theorem: they
are data needed to formulate the theorem. Supplying the channel properties or
the normalization as assumptions would yield only a conditional statement and
would not formalize the implication asserted by the paper.

\section{Formal boundary}

The source-faithful local contraction
\eqref{eq:cpsv16_peps_boundary_rfp_specification_local_tensor} is complete.
The implication from the PEPS renormalization diagram to the boundary MPO
renormalization equations remains an open statement-level problem. A faithful
completion must first choose an algebraic PEPS blocking relation and boundary
fixed-point construction that determine the two channels, and must then derive
their channel properties and
Eqs.~\eqref{eq:cpsv16_peps_boundary_rfp_specification_coarse_graining}--
\eqref{eq:cpsv16_peps_boundary_rfp_specification_refinement}. No additional
complete-positivity, trace-preservation, fixed-point normalization, or
extension hypothesis may be inserted into the paper-facing implication.

This boundary is tracked in \ghissue{7371}.

\bibliographystyle{alpha}
\bibliography{references}

\end{document}
